\documentclass[11pt,a4paper]{article}

%% PACHETE MATEMATICE
\usepackage{amsmath,amssymb,amsfonts}
\usepackage{amsthm}
\usepackage{mathpazo}
\usepackage{eulervm}
\usepackage{enumerate}
\usepackage{color}
\usepackage{graphicx}
\usepackage[all]{xy}
\usepackage{xcolor}
\usepackage{framed}
\usepackage[unicode]{hyperref}
\setcounter{MaxMatrixCols}{20}
\usepackage{multicol}
%\usepackage[romanian]{babel}


%% ASPECT
\linespread{1.05}
\usepackage[T1]{fontenc}
\usepackage[utf8x]{inputenc}
\usepackage[margin=2cm]{geometry}
% \usepackage[color]{showkeys}
\usepackage{anyfontsize}
\usepackage{titlesec}
\titleformat*{\section}{\large\bfseries}

\usepackage{fancyhdr}
\pagestyle{fancy}
\rhead{\small \color{gray} Notițe de Adrian Manea}
\lhead{\small \color{gray} M1-ACS, 2017---2018, M. Olteanu}


\usepackage[autostyle = false, style = german]{csquotes}
\usepackage[romanian]{babel}
\newcommand\qq{\enquote}

%% COMENZILE MELE
\renewcommand*{\proofname}{Dem.:}
\newcommand\bb{\mathbb}
\newcommand\dr{\mathrm}
\newcommand\kal{\mathcal}
\newcommand\fk{\mathfrak}
\newcommand\op{\oplus}
\newcommand\ot{\otimes}
\newcommand\ds{\displaystyle}
\newcommand\id{\indent}
\newcommand\nid{\noindent}
\newcommand\seq{\subseteq}
\newcommand\sneq{\subsetneq}
\newcommand\speq{\supseteq}
\newcommand\spneq{\supsetneq}
\newcommand\rar{\rightarrow}
\newcommand\Rar{\Rightarrow}
\newcommand\xrar{\xrightarrow}
\newcommand\lar{\leftarrow}
\newcommand\Lar{\Leftarrow}
\newcommand\xlar{\xleftarrow}
\newcommand\lrar{\leftrightarrow}
\newcommand\Lrar{\Leftrightarrow}
\newcommand\ideal{\trianglelefteq}
\newcommand\ai{\text{ a.\^{i}. }}
\newcommand\sk{(\textit{Schi\c{t}\u{a}}) }
\newcommand\rhup{\rightharpoonup}
\newcommand\rhdn{\rightharpoondown}
\newcommand\lhup{\leftharpoonup}
\newcommand\lhdn{\leftharpoondown}
\newcommand\vid{\emptyset}
\newcommand\wh{\widehat}
\newcommand\ol{\overline}
\newcommand\NN{\mathbb{N}}
\newcommand\RR{\mathbb{R}}
\newcommand\ZZ{\mathbb{Z}}
\newcommand\QQ{\mathbb{Q}}
\newcommand\CC{\mathbb{C}}

%% STILURILE MELE
\newtheoremstyle{dotless}{}{}{\itshape}{}{\bfseries}{:}{ }{}

\theoremstyle{dotless}
\newtheorem{theorem}{Teorem\u{a}}[section]
\newtheorem{proposition}{Propozi\c{t}ie}[section]
\newtheorem{lemma}{Lem\u{a}}[section]
\newtheorem{corollary}{Corolar}[section]
\newtheorem{exercise}{Exerci\c{t}iu}

\newtheoremstyle{dotlessdr}{}{}{}{}{\bfseries}{:}{ }{}
\theoremstyle{dotlessdr}
\newtheorem{example}{Exemplu}[section]
\newtheorem{definition}{Defini\c{t}ie}[section]
\newtheorem{remark}{Observa\c{t}ie}[section]




\begin{document}

\begin{center}
  {\large\textbf{
      Seminar 9 \\
      Extreme cu legături. Integrale improprii
    }}
\end{center}

\vspace{1cm}

\section{Extreme condiționate}
\label{sec:extreme-cond}

Atunci cînd domeniul de definiție al unei funcții de mai multe variabile conține,
la rîndul său anumite ecuații (numite, generic, \emph{legături}, problemele de
extrem se studiază folosind \textbf{metoda multiplicatorilor Lagrange}. Ea se
bazează pe următoarele concepte:

\begin{definition}\label{def:fc-scop}
  Fie $f(\pmb{x}, \pmb{y})$, cu $\pmb{x} \in \RR^n, \pmb{y} \in \RR^m$,
  iar $f : U \to \RR$ o funcție de $n + m$ variabile reale, cu valori reale, de
  clasă $\kal{C}^1$ pe un deschis $U \seq \RR^{n+m}$.

  Ea se numește \emph{funcție-scop (obiectiv)}. Presupunem că există $m$
  legături între variabilele $\pmb{x}, \pmb{y}$, adică $m$ relații de forma:
  \[
    g_1(\pmb{x}, \pmb{y}) = 0, \dots, g_m(\pmb{x}, \pmb{y}) = 0, \quad %
    g_i : U \to \RR,
  \]
  fiecare legătură fiind o funcție de clasă $\kal{C}^1(U)$.

  Fie $M = \{(x, y) \in U \mid g_i(\pmb{x}, \pmb{y}) = 0, 1 \leq i \leq m\}$
  mulțimea punctelor din $U$ care verifică legăturile.

  Se numește \emph{punct de extrem local al funcției $f$ cu legăturile $g_i$}
  orice punct $(\pmb{x}_0, \pmb{y}_0) \in M$, pentru care există o vecinătate
  $W \seq U$, astfel încît diferența $f(\pmb{x}, \pmb{y}) - f(\pmb{x}_0, \pmb{y}_0)$
  să aibă semn constant pentru orice $(\pmb{x}, \pmb{y}) \in M \cap W$.
\end{definition}

Teorema pe care se bazează metoda multiplicatorilor lui Lagrange este:

\begin{theorem}[J. L. Lagrange]\label{thm:lagr-mult}
  Cu notațiile și contextul de mai sus, presupunem că $(\pmb{x}_0, \pmb{y}_0)$
  este un punct de extrem local al lui $f$, cu legăturile de mai sus și că
  \[
    \det J_g(\pmb{x}_0, \pmb{y}_0) \stackrel{\text{not.}}{=\joinrel=\joinrel=} %
    \frac{D(g_1, \dots, g_m)}{D(y_1, \dots, y_m)} (\pmb{x}_0, \pmb{y}_0) \neq 0.
  \]

  Atunci există $m$ numere reale $\lambda_1, \dots, \lambda_m$, numite
  \emph{multiplicatori Lagrange} astfel încît, dacă definim funcția:
  \[
    F = f + \lambda_1 g_1 + \dots + \lambda_m g_m,
  \]
  punctul $(\pmb{x}_0, \pmb{y}_0)$ să verifice în mod necesar sistemul de
  $n + 2m$ ecuații:
  \[
    \frac{\partial F}{\partial x_j} = \frac{\partial F}{\partial y_k} = 0, %
    g_l = 0, \quad \forall 1 \leq j \leq n, 1 \leq k \leq m, 1 \leq l \leq m,
  \]
  cu $n + 2m$ necunoscute, $\pmb{\lambda}, \pmb{x}, \pmb{y}$.
\end{theorem}

\textbf{Exemplu:} Să se determine extremele locale ale funcției $f(x,y,z) = xyz$,
cu legătura $x + y + z = 1$.

\emph{Soluție:} Folosim metoda multiplicatorilor Lagrange. Definim funcțiile:
\begin{align*}
  g(x, y, z) &= x + y + z - 1 \\
  F(x, y, z) &= f + \lambda g = xyz + \lambda(x + y + z - 1).
\end{align*}

Extremele cerute verifică sistemul:
\[
  \begin{cases}
    \frac{\partial F}{\partial x} &= 0 \\
    \frac{\partial F}{\partial y} &= 0 \\
    \frac{\partial F}{\partial z} &= 0 \\
    g &= 0
  \end{cases} \Rightarrow
  \begin{cases}
    yz + \lambda &= 0 \\
    xz + \lambda &= 0 \\
    xy + \lambda &= 0 \\
    x + y + z - 1 &= 0
  \end{cases}
\]

Soluția sistemului este dată de:
\begin{align*}
  \lambda_1 &= -\frac{1}{9} \Rightarrow %
              (x_1, y_1, z_1) = \big(\frac{1}{3}, \frac{1}{3}, \frac{1}{3}\big) \\
  \lambda_2 &= 0 \Rightarrow
              \begin{cases}
                (x_2, y_2, z_2) &= (1, 0, 0) \\
                (x_3, y_3, z_3) &= (0, 1, 0) \\
                (x_4, y_4, z_4) &= (0, 0, 1)
              \end{cases}
\end{align*}

În continuare, studiem natura punctelor de extrem pentru funcția $F$, fie cu
matricea hessiană, fie cu diferențiala totală de ordin 2.

Găsim că $(x_1, y_1, z_1)$ este maxim local, iar celelalte puncte nu sînt
de extrem.

%%%%%%%%%%%%%%%%%%%%%%%%%%%%%%%%%%%%%%%%%%%%%%%%%%%%%%%%%%%%%%%%%%%%%%

\section{Integrale improprii. Criterii de convergență}
\label{sec:int-imp}

Integralele improprii reprezintă cazul în care funcția care se integrează
nu este mărginită la cel puțin unul dintre capetele domeniului de integrare.

Fie $a, b \in \RR$ și $f : [a, b) \to \RR$ o funcție \emph{local integrabilă}
(i.e.\ integrabilă pe orice interval compact $[u, v] \seq [a, b)$).

Integrala improprie (în $b$) $\ds\int_a^b f(x) dx$ se numește \emph{convergentă}
dacă limita:
\[
  \lim_{t \to b} \int_a^t f(x) dx
\]
există și este finită. Valoarea limitei este valoarea integralei. În caz contrar,
integrala se numește \emph{divergentă}.

Dacă $f : [a, \infty) \to \RR$ este local integrabilă, atunci integrala
improprie (la $\infty$) $\ds\int_a^\infty f(x) dx$ se numește \emph{convergentă}
dacă limita:
\[
  \lim_{t \to \infty} \int_a^t f(x) dx
\]
există și este finită. Valoarea limitei este egală cu valoarea integralei.

Integrala improprie $\ds\int_a^b f(x) dx$ se numește \emph{absolut convergentă}
dacă integrala $\ds\int_a^b |f(x)| dx$ este convergentă.

Criteriile de convergență pentru integralele improprii sînt foarte asemănătoare
cu cele pentru serii (amintiți-vă, că, de fapt, integralele definite se construiesc
cu ajutorul sumelor infinite, adică serii, v.\ \emph{sumele Riemann}).

Așadar, avem:

\textbf{Criteriul lui Cauchy} (general): Fie $f:[a, b) \to \RR$ local integrabilă.
Atunci integrala $\ds\int_a^b f(t) dt$ este convergentă dacă și numai dacă:
\[
  \forall \varepsilon > 0, \exists b_\varepsilon \in [a, b) \ai %
  \forall x, y \in (b_\varepsilon, b), \big| \int_x^y f(t) dt \big| < \varepsilon.
\]

\vspace{1cm}

\textbf{Criteriul de comparație} (\qq{termen cu termen}): Fie $f, g : [a, b) \to \RR$
astfel încît $0 \leq f \leq g$.
\begin{itemize}
\item Dacă $\ds\int_a^b g(x) dx$ este convergentă, atunci și integrala
  $\ds\int_a^b f(x) dx$ este convergentă;
\item Dacă integrala $\ds\int_a^b f(x) dx$ este divergentă, atunci și integrala
  $\ds\int_a^b g(x) dx$ este divergentă.
\end{itemize}

\vspace{1cm}

\textbf{Criteriul de comparație la limită}: Fie $f, g : [a, b) \to [0, \infty)$,
astfel încît să existe limita:
\[
  \ell = \lim_{x \to b} \frac{f(x)}{g(x)}.
\]
\begin{itemize}
\item Dacă $\ell \in [0, \infty)$, iar $\ds\int_a^b g(x) dx$ este convergentă,
  atunci $\ds\int_a^b f(x) dx$ este convergentă;
\item Dacă $\ell \in (0, \infty)$ sau $\ell = \infty$, iar $\ds\int_a^b g(x) dx$
  este divergentă, atunci și $\ds\int_a^b f(x) dx$ este divergentă.
\end{itemize}

\vspace{1cm}

\textbf{Criteriul de comparație cu $\frac{1}{x^\alpha}$}:
Fie $a \in \RR$ și $f : [a, \infty) \to [0, \infty)$ local integrabilă, astfel
încît să existe:
\[
  \ell = \lim_{x \to \infty} x^\alpha f(x).
\]
\begin{itemize}
\item Dacă $\alpha > 1$ și $0 \leq \ell < \infty$, atunci $\ds\int_a^\infty f(x) dx$
  este convergentă;
\item Dacă $\alpha \leq 1$, iar $0 < \ell \leq \infty$, atunci $\ds\int_a^\infty f(x) dx$
  este divergentă.
\end{itemize}

\vspace{1cm}

\textbf{Criteriul de comparație cu $\frac{1}{(b - x)^\alpha}$}: Fie $a < b$ și
$f : [a, b) \to [0, \infty)$, local integrabilă, astfel încît să existe:
\[
  \ell = \lim_{x \to b} (b - x)^\alpha f(x).
\]
\begin{itemize}
\item Dacă $\alpha < 1$ și $0 \leq \ell < \infty$, atunci $\ds\int_a^b f(x) dx$
  este convergentă;
\item Dacă $\alpha \geq 1$ și $0 < \ell \leq \infty$, atunci $\ds\int_a^b f(x) dx$
  este divergentă.
\end{itemize}

\vspace{1cm}

\textbf{Criteriul lui Abel}: Fie $f, g : [a, \infty) \to \RR$, cu proprietățile:
\begin{itemize}
\item $f$ este de clasă $\kal{C}^1$, $\ds\lim_{x \to \infty} f(x) = 0$, iar
  $\ds\int_a^\infty f'(x) dx$ este absolut convergentă;
\item $g$ este continuă, iar $G(x) = \ds\int_a^x f(t) dt$ este mărginită pe
  $[a, \infty)$.
\end{itemize}
Atunci integrala $\ds\int_a^\infty f(x) g(x) dx$ este convergentă.

%%%%%%%%%%%%%%%%%%%%%%%%%%%%%%%%%%%%%%%%%%%%%%%%%%%%%%%%%%%%%%%%%%%%%%

\section{Integrale cu parametri}
\label{sec:int-param}

Fie $A \neq \vid$ și $[a, b] \seq \RR$ un interval compact. Considerăm funcția
$f : [a, b] \times A \to \RR$, astfel încît, pentru orice $y \in A$, funcția
$[a, b] \ni x \mapsto f(x, y) \in \RR$ să fie integrabilă Riemann.

Funcția definită prin:
\[
  F : A \to \RR, \quad F(y) = \int_a^b f(x, y) dx
\]
se numește \emph{integrală cu parametru}.

Proprietățile pe care le vom utiliza sînt cele de mai jos.

\textbf{Continuitate}: Dacă $f : [a, b] \times A \to \RR$ este continuă, atunci
integrala cu parametru $F(y)$ definită mai sus este funcție continuă.

\vspace{1cm}

\textbf{Formula de derivare (Leibniz)}: Fie $f : [a, b] \times (c, d) \to \RR$
o funcție continuă, astfel încît derivata parțială $\frac{\partial f}{\partial y}$
există și este continuă pe $[a, b] \times (c, d)$. Atunci integrala cu parametru
$F(y)$ definită mai sus este derivabilă și are loc:
\[
  F'(y) = \int_a^b \frac{\partial f}{\partial y}(x, y) dx, \forall y \in (c, d).
\]

\vspace{1cm}

\textbf{Formula generală de derivare:} Dacă $f : [a, b] \times (c, d) \to \RR$
este o funcție continuă, astfel încît derivata parțială $\frac{\partial f}{\partial y}$
să existe și să fie continuă pe $[a, b] \times (c, d)$, definim
$\varphi, \psi : (c, d) \to [a, b)$ două funcții de clasă $\kal{C}^1$.

Atunci funcția $G(y) = \int_{\varphi(y)}^{\psi(y)} f(x, y) dx$ este derivabilă
și are loc formula de derivare:
\[
  G'(y) = \int_{\varphi(y)}^{\psi(y)} \frac{\partial f}{\partial y}(x, y) dx + %
  f(\varphi(y), y) \psi'(y) - f(\varphi(y), y) \varphi'(y), \forall y \in (c, d).
\]

\vspace{1cm}

\textbf{Schimbarea ordinii de integrare:} Dacă $f: [a, b] \times [c, d] \to \RR$
este o funcție continuă, atunci are loc:
\[
  \int_a^b \Big( \int_c^d f(x, y) dy \Big) dx = %
  \int_c^d \Big( \int_a^b f(x, y) dx \Big) dy.
\]

%%%%%%%%%%%%%%%%%%%%%%%%%%%%%%%%%%%%%%%%%%%%%%%%%%%%%%%%%%%%%%%%%%%%%%

\section{Integrale improprii cu parametri}
\label{sec:int-imp-par}

Putem considera acum integrale improprii, definite cu parametri, astfel.
Luăm o funcție $f : [a, b) \times A \to \RR$, astfel încît pentru orice $y \in A$,
aplicația $[a, b) \ni x \mapsto f(x, y) \in \RR$ este local integrabilă și integrala
$\ds\int_a^b f(x, y) dx$ converge. Atunci putem defini funcția:
\[
  F(x,y) = \int_a^b f(x, y) dx,
\]
care se numește \emph{integrală improprie cu parametru}.

\begin{definition}\label{def:uc}
  Integrala $F(x, y)$ de mai sus se numește \emph{uniform convergentă} (UC) (în raport cu $y$)
  pe mulțimea $A$ dacă:
  \[
    \forall \varepsilon > 0, \exists b_\varepsilon \in (a, b) \ai %
    \Big| \int_t^b f(x, y) dx \Big| < \varepsilon, \forall t \in (b_\varepsilon, b), %
    \forall y \in A.
  \]
\end{definition}

Pentru aceste integrale, se pot adapta proprietățile integralelor cu parametri
din secțiunea anterioară:

\textbf{Continuitate}: Dacă $f : [a, b) \times A \to \RR$ este continuă, iar
integrala $\ds\int_a^b f(x, y) dx$ este UC pe $A$, atunci funcția $F(x,y)$
definită mai sus este continuă.

\vspace{1cm}

\textbf{Derivare}: Fie $f: [a, b) \times (c, d) \to \RR$ o funcție continuă, astfel
încît derivata parțială $\frac{\partial f}{\partial y}$ există și este continuă pe
$[a, b) \times (c, d)$ și pentru orice $y \in (c, d)$ fixat, integrala
$F(y) = \ds\int_a^b f(x,y) dx$ este convergentă.

Dacă integrala $\ds\int_a^b \frac{\partial f}{\partial y}(x, y) dx$ este UC
pe $(c, d)$, atunci integrala improprie cu parametru $F(y)$ de mai sus este
derivabilă și are loc:
\[
  F'(y) = \int_a^b \frac{\partial f}{\partial y}(x, y) dx, \forall y \in (c, d).
\]

\vspace{1cm}

\textbf{Schimbarea ordinii de integrare}: Dacă $f : [a, b) \times [c, d] \to \RR$
este continuă și integrala $F(y) = \ds\int_a^b f(x, y) dx$ este UC pe $(c, d)$,
atunci are loc:
\[
  \int_c^d \Big( \int_a^b f(x, y) dx \Big) dy = %
  \int_a^b \Big( \int_c^d f(x, y) dy \Big) dx.
\]

\vspace{1cm}

\textbf{Criteriul de comparație pentru UC}: Fie $f: [a, b) \times A \to \RR$ o
funcție cu proprietatea că, pentru orice $y \in A$, aplicația
$[a, b) \ni x \mapsto f(x, y) \in \RR$ este local integrabilă. Fie $g : [a, b) \to \RR$,
astfel încît $|f(x, y)| \leq g(x), \forall x \in [a, b), \forall y \in A$.

Dacă integrala $\ds\int_a^b g(x) dx$ este convergentă, atunci integrala
$\ds\int_a^b f(x, y) dx$ este UC.

%%%%%%%%%%%%%%%%%%%%%%%%%%%%%%%%%%%%%%%%%%%%%%%%%%%%%%%%%%%%%%%%%%%%%%

\subsection{Funcțiile lui Euler}
\label{subsec:f-euler}

Următoarele integrale improprii cu parametri se numesc \emph{funcțiile lui Euler}:
\begin{align*}
  \Gamma(\alpha) &= \int_0^\infty x^{\alpha - 1} e^{-x} dx, \alpha > 0 \\
  B(p, q) &= \int_0^1 x^{p-1} (1 - x)^{q - 1} dx, p > 0, q > 0.
\end{align*}

Proprietățile lor, pe care le vom utiliza în calcule, sînt:
\begin{itemize}
\item $B(p, q) = B(q, p)$;
\item $B(p, q) = \frac{\Gamma(p)\Gamma(q)}{\Gamma(p + q)}$;
\item $B(p, q) = \ds\int_0^\infty \frac{y^{p-1}}{(1 + y)^{p + q}} dy$;
\item $\Gamma(1) = 1$;
\item $\Gamma(\alpha + 1) = \alpha \cdot \Gamma(\alpha)$;
\item $\Gamma(n) = (n - 1)!, \forall n \in \NN$;
\item $\Gamma\big( \frac{1}{2} \big) = \sqrt{\pi}$;
\item $\Gamma\big( n + \frac{1}{2} \big) = (2n - 1)!! \cdot 2^{-n} \cdot \sqrt{\pi}, %
  \forall n \in \NN$;
\item $\Gamma(\alpha) \Gamma(1 - \alpha) = \frac{\pi}{\sin(\alpha \pi)}, %
  \forall \alpha \in (0, 1)$.
\end{itemize}

%%%%%%%%%%%%%%%%%%%%%%%%%%%%%%%%%%%%%%%%%%%%%%%%%%%%%%%%%%%%%%%%%%%%%%

\section{Exerciții}
\label{sec:exc}

1. Folosind criteriile de comparație, să se studieze natura integralelor improprii:
\begin{enumerate}[(a)]
\item $\ds\int_1^\infty \frac{dx}{\sqrt{x^2 + 1}}$ \hfill(D);
\item $\ds\int_0^1 \frac{x^2}{\sqrt{1 - x^2}} dx$ \hfill(C);
\item $\ds\int_0^1 \frac{\sin x}{1 - x^2} dx$ \hfill(D);
\item $\ds\int_1^\infty \frac{x}{\sqrt{x^3 - 1}} dx$ \hfill(C $x = 1$, D
  $x \to \infty \Rightarrow$ D);
\item $\ds\int_1^\infty \frac{\ln x}{\sqrt{x^3 - 1}} dx$ \hfill(C);
\item $\ds\int_1^\infty \frac{dx}{x \sqrt{x} - 1} dx$ \hfill(C $x \to \infty$,
  D $x = 1 \Rightarrow$ D);
\item $\ds\int_1^\infty \frac{dx}{x(\ln x)^\alpha}, \alpha > 0$ \hfill(D);
\end{enumerate}

\vspace{1cm}

2. Să se arate că integrala $\ds\int_0^\infty \frac{\sin x}{x} dx$ este convergentă,
dar nu este absolut convergentă.

\emph{Indicație:} Pentru $x \to \infty$, convergența rezultă din criteriul
lui Abel. Pentru $x = 0$, putem prelungi funcția prin continuitate, deoarece
limita sa este finită.

Pentru AC, se aplică criteriul de comparație.

\vspace{1cm}

3. Să se calculeze integralele, folosind derivarea sub integrală:
\begin{enumerate}[(a)]
\item $I(m) = \ds\int_0^{\frac{\pi}{2}} \ln(\cos^2 x + m^2 \sin^2 x) dx, m > 0$;
\item $I(a) = \ds\int_0^{\frac{\pi}{2}} \frac{\ln(1 + a \cos x)}{\cos x} dx, |a| < 1$;
\end{enumerate}

\emph{Indicații:}
\begin{enumerate}[(a)]
\item Derivăm în raport cu $m$ sub integrală, apoi integrăm cu schimbarea de variabilă
  $\tan x = t$;
\item Derivăm în raport cu $a$ sub integrală și apoi integrăm cu schimbarea de variabilă
  $\tan \frac{x}{2} = t$.
\end{enumerate}

\vspace{1cm}

4. Să se calculeze, folosind funcțiile $\Gamma$ și $B$, integralele:
\begin{enumerate}[(a)]
\item $\ds\int_0^\infty e^{-x^p} dx, p > 0$;
\item $\ds\int_0^{\infty} \frac{x^{\frac{1}{4}}}{(x + 1)^2} dx$;
\item $\ds\int_0^\infty \frac{dx}{x^3 + 1} dx$;
\item $\ds\int_0^{\frac{\pi}{2}} \sin^p x \cos^q x dx, p > -1, q > - 1$;
\item $\ds\int_0^1 x^{p+1}(1-x^m)^{q - 1} dx, p, q, m > 0$;
\item $\ds\int_0^\infty x^p e^{-x^q} dx, p > -1, q > 0$;
\item $\ds\int_0^1 \ln^p x^{-1} dx, p > -1$;
\item $\ds\int_0^1 \frac{dx}{(1 - x^n)^{\frac{1}{n}}}, n \in \NN$.
\end{enumerate}

\emph{Indicații:}
\begin{enumerate}[(a)]
\item Facem schimbarea de variabilă $x^p = y$ și obținem $\Gamma\Big(\frac{1}{p} + 1\Big)$;
\item Folosind proprietățile, obținem $B \Big( \frac{5}{4}, \frac{3}{4} \Big)$, pe care
  îl scriem în funcție de $\Gamma$. Răspuns: $\frac{\pi \sqrt{2}}{4}$;
\item Facem schimbarea de variabilă $x^3 = y$ și găsim
  $\frac{1}{3} B \Big( \frac{1}{3}, \frac{2}{3} \Big)$;
\item Facem schimbarea de variabilă $\sin^2 x = y$ și scriem în funcție de $B$;
\item $x^m = y$ și scriem în funcție de $B$
\item $x^q = y$ și scriem în funcție de $\Gamma$;
\item $\ln x^{-1} = y$ și scriem în funcție de $\Gamma$;
\item $x^n = y$ și scriem în funcție de $B$.
\end{enumerate}

\vspace{1cm}

5. Să se determine extremele funcțiilor $f$, cu legătura $g$ în cazurile:
\begin{enumerate}[(a)]
\item $f(x, y) = x^2 + y^2 - 6x - 2y + 1, g(x, y) = x + y - 2$;
\item $f(x, y) = x^2 + y^2 - 6x - 2y + 1, g(x, y) = x^2 + y^2 - 2x - 2y + 1$;
\item $f(x, y) = 3x + 4y, g(x, y) = x^2 + y^2 + 25$.
\end{enumerate}

\vspace{1cm}

6. Să se găsească punctul din planul $2x + y - z = 5$, situat la distanță minimă
față de origine.

\vspace{1cm}

7. Să se determine punctele cele mai depărtate de origine care se află pe
suprafața de ecuație $4x^2 + y^2 + z^2 - 8x - 4y + 4 = 0$.

\vspace{1cm}

8. Să se determine valorile extreme ale funcției:
\[
  f: \RR^3 \to \RR, \quad f(x, y, z) = 2x^2 + y^2 + 3z^2,
\]
pe mulțimea $\{(x, y, z) \in \RR^3 \mid x^2 + y^2 + z^2 = 1 \}$.

\vspace{1cm}

9. Să se determine valorile extreme ale produsului $xy$, cînd $x$ și $y$ sînt
coordonatele unui punct de pe elipsa de ecuație $x^2 + 2y^2 = 1$.
\end{document}
